\chapter{Fundamentos del aprendizaje automático y modelos clásicos}
\section{Apuntes del Tema 1: Foundations of Machine Learning — Classical Models \& Clustering}

\subsection{1. Conceptos Fundamentales y Paradigmas de Aprendizaje}

\subsubsection{1.1 Trazabilidad y Definición}

\begin{itemize}
	\item \textbf{Machine Learning:} Disciplina enfocada en la creación de sistemas capaces de realizar tareas mediante el reconocimiento de patrones y la toma de decisiones basadas en datos de experiencia, evaluando el desempeño mediante la capacidad de \textbf{generalización} hacia nuevos datos no vistos y previniendo el \emph{overfitting}.
	
	\item \textbf{Estructura de Datos Básica:}
	\begin{itemize}
		\item \textbf{Dataset:} Conjunto de observaciones $\mathcal{D} = \{(x^{(i)}, y^{(i)})\}_{i=1}^{n}$.
		\item \textbf{Instancia ($x$):} Vector de características de dimensión $d$, donde $x \in \mathbb{R}^d$.
		\item \textbf{Etiqueta/Target ($y$):} Valor objetivo o clase a predecir.
		\item \textbf{Modelo Paramétrico:} Función de mapeo $f(x; \theta) \rightarrow \hat{y}$ parametrizada por pesos y sesgos $\theta = (w, b)$.
	\end{itemize}
\end{itemize}

\subsubsection{1.2 Paradigmas de Aprendizaje}

\begin{itemize}
	\item \textbf{Supervisado:} Aprendizaje mediante pares entrada-salida etiquetados $(x, y)$. Se divide en:
	\begin{itemize}
		\item \textbf{Regresión:} El target $y$ es un valor continuo.
		\item \textbf{Clasificación:} El target $y$ representa una categoría discreta.
	\end{itemize}
	
	\item \textbf{No Supervisado:} Descubrimiento de estructuras implícitas o agrupamientos en datos sin etiquetas $x$ (ej. \emph{Clustering}, reducción de dimensionalidad).
	
	\item \textbf{Otros Paradigmas Relevantes:}
	\begin{itemize}
		\item \emph{Semi-supervised}: Combina una pequeña cantidad de datos etiquetados con una gran cantidad de datos no etiquetados para mejorar la generalización.
		\item \emph{Self-supervised}: Aprendizaje donde el propio modelo genera las etiquetas a partir de los datos mediante tareas pretexto (\emph{pretext tasks}) para capturar representaciones significativas.
		\item \emph{Transfer Learning}: Reutilización de modelos o backbones previamente entrenados en otra tarea o dominio para reducir la necesidad de datos.
		\item \emph{Active Learning}: Estrategia en la que el modelo interactúa seleccionando los puntos de datos más informativos para que un humano los etiquete, reduciendo costes.
		\item \emph{Online/Incremental Learning}: Actualización continua del modelo en tiempo real a medida que fluyen nuevos datos, útil para entornos dinámicos.
		\item \emph{Federated Learning}: Entrenamiento colaborativo de un modelo global de forma distribuida en dispositivos locales sin compartir los datos privados de los usuarios.
	\end{itemize}
\end{itemize}

\subsection{2. k-Nearest Neighbors (kNN)}

\subsubsection{2.1 Principio e Intuición}

\begin{itemize}
	\item \textbf{Algoritmo Lazy Learner:} No requiere una fase de entrenamiento o ajuste de parámetros explícito ($f(x)$ se evalúa en tiempo de consulta).
	
	\item \textbf{Regla de Clasificación:} Asigna la etiqueta a una muestra de prueba $x_{\text{query}}$ según el voto mayoritario (o promedio ponderado por la distancia inversa) de sus $k$ vecinos más cercanos en el espacio métrico.
\end{itemize}

\subsubsection{2.2 Métricas de Distancia}

\begin{itemize}
	\item \textbf{Distancia Euclídea ($L_2$):}
	$$\text{dist}_{\text{Euclidean}}(x, x') = \sqrt{\sum_{j=1}^{d} (x_j - x_j')^2}$$
	
	\item \textbf{Distancia Manhattan ($L_1$):}
	$$\text{dist}_{\text{Manhattan}}(x, x') = \sum_{j=1}^{d} |x_j - x_j'|$$
	
	\item \textbf{Minkowski ($L_p$):} 
	$$\text{dist}_{L^p}(x, x') = \left( \sum_{j=1}^{d} |x_j - x_j'|^p \right)^{1/p}$$
\end{itemize}

\subsubsection{2.3 Selección de $k$ y Efectos de Escala}

\begin{itemize}
	\item \textbf{Efecto de $k$:}
	\begin{itemize}
		\item $k=1$ (Bajo): Frontera extremadamente flexible, muy sensible al ruido y propenso a sobreajuste (\emph{overfitting}).
		\item $k$ elevado: Fronteras muy suaves, riesgo de subajuste (\emph{underfitting}) al mezclas vecindarios.
	\end{itemize}
	
	\item \textbf{Necesidad de Normalización:} Dado que kNN depende de distancias geométricas, las variables con rangos numéricos elevados dominan el cálculo. Es indispensable estandarizar las características ($x' = \frac{x - \mu}{\sigma}$) antes de calcular las distancias.
	
	\item \textbf{La Maldición de la Dimensionalidad (\emph{Curse of Dimensionality}):} A medida que la dimensión $d$ del espacio crece, el volumen aumenta exponencialmente, volviendo los datos muy dispersos y reduciendo la capacidad discriminativa de las distancias geométricas entre los puntos.
\end{itemize}

\subsection{3. Naive Bayes}

\subsubsection{3.1 Fundamento Probabilístico Generativo}

\begin{itemize}
	\item \textbf{Modelo Generativo:} Modela la probabilidad de la clase a priori $P(y)$ y la verosimilitud de los datos $P(x|y)$ para estimar la probabilidad a posteriori $P(y|x)$ aplicando el Teorema de Bayes:
	$$P(y|x) = \frac{P(x|y)P(y)}{P(x)} \propto P(y) \prod_{j=1}^{d} P(x_j|y)$$
	
	\item \textbf{Regla MAP (Maximum A Posteriori):} Selección de la clase que maximiza el producto: 
	$$\hat{y} = \arg\max_{y} P(y) \prod_{j=1}^{d} P(x_j|y)$$
	
	\item \textbf{Asunción "Naive" (Ingenua):} Presupone la \textbf{independencia condicional} de todas las características $x_j$ dada la clase $y$.
\end{itemize}

\subsubsection{3.2 Variantes Principales de Naive Bayes}

\begin{table}[htbp]
	\centering
	\small
	\begin{tabular}{|p{3cm}|p{3.5cm}|p{5.5cm}|}
		\hline
		\textbf{Variante} & \textbf{Tipo de Variable ($x_j$)} & \textbf{Estimación de Verosimilitud $P(x_j|y)$} \\ \hline
		\textbf{Categorical NB} & Categóricas/Nominales discretas & $P(x_j = v \mid y) = \frac{\text{count}(x_j=v, y) + \alpha}{\text{count}(y) + \alpha \cdot K_j}$ \\ \hline
		\textbf{Multinomial NB} & Conteos discretos (frecuencias de texto) & $P(x_j \mid y) = \frac{\text{count}(x_j \text{ en clase } y) + \alpha}{\sum_{k} \text{count}(x_k \text{ en clase } y) + \alpha \cdot V}$ \\ \hline
		\textbf{Gaussian NB} & Características continuas & $P(x_j \mid y) = \frac{1}{\sqrt{2\pi \sigma_{j,y}^2}} \exp\left(-\frac{(x_j - \mu_{j,y})^2}{2\sigma_{j,y}^2}\right)$ \\ \hline
	\end{tabular}
\end{table}

\begin{itemize}
	\item \textbf{Laplace Smoothing ($\alpha=1$):} Corrige el problema de la \textbf{probabilidad cero} sumando el parámetro $\alpha$ en el numerador y $\alpha \cdot K_j$ (o $\alpha \cdot V$) en el denominador.
\end{itemize}

\subsection{4. Árboles de Decisión (Decision Trees)}

\subsubsection{4.1 Concepto y Particionado Recursivo}

\begin{itemize}
	\item Modelos no paramétricos e interpretables para clasificación y regresión que dividen recursivamente el espacio de características mediante reglas ortogonales a los ejes.
\end{itemize}

\subsubsection{4.2 Comparativa de Algoritmos y Criterios de División}

\begin{enumerate}
	\item \textbf{ID3 (Iterative Dichotomiser 3):}
	\begin{itemize}
		\item Solo admite atributos categóricos y no realiza poda.
		\item Utiliza la \textbf{Entropía} $H(S)$ y la \textbf{Ganancia de Información} $IG(S, A)$:
		$$H(S) = - \sum_{i=1}^{c} p_i \log_2(p_i)$$
		$$IG(S, A) = H(S) - \sum_{v \in \text{Valores}(A)} \frac{|S_v|}{|S|} H(S_v)$$
	\end{itemize}
	
	\item \textbf{C4.5:}
	\begin{itemize}
		\item Soporta variables continuas mediante la evaluación de umbrales binarios de corte $t$.
		\item Sustituye el $IG$ por el \textbf{Gain Ratio} para corregir el sesgo de ID3 hacia atributos con excesivos valores únicos:
		$$\text{GainRatio}(S, A) = \frac{IG(S, A)}{\text{SplitInfo}(A)}$$
		$$\text{SplitInfo}(A) = -\sum_{i=1}^{k} \frac{|S_i|}{|S|} \log_2\left(\frac{|S_i|}{|S|}\right)$$
		\item Aplica poda pesimista posterior al entrenamiento.
	\end{itemize}
	
	\item \textbf{CART (Classification and Regression Trees):}
	\begin{itemize}
		\item Realiza exclusivamente divisiones binarias.
		\item Minimiza la \textbf{Impureza de Gini} en clasificación:
		$$G(S) = 1 - \sum_{i=1}^{c} p_i^2$$
		\item En regresión, minimiza la varianza/suma de errores cuadráticos (SSE).
	\end{itemize}
\end{enumerate}

\subsection{5. Clustering (Agrupamiento No Supervisado)}

\subsubsection{5.1 k-means}

\begin{itemize}
	\item \textbf{Objetivo:} Particionar $X$ en $c$ clusters disjuntos $\omega_1, \dots, \omega_c$ minimizando la varianza dentro de cada cluster (suma de distancias al cuadrado al centroide $\mu_i$).
	
	\item \textbf{Paso a Paso del Algoritmo:}
	\begin{enumerate}
		\item Inicialización aleatoria de $k$ centroides.
		\item Asignación de cada punto al centroide más cercano.
		\item Reordenamiento y recalculo del punto medio/centroide de cada cluster.
		\item Repetición iterativa hasta convergencia.
	\end{enumerate}
	
	\item \textbf{Selección de $k$:} Se evalúa empíricamente mediante el \textbf{Elbow Method} y métricas como el \textbf{Silhouette Score} o el \textbf{Davies–Bouldin Score}.
\end{itemize}

\subsubsection{5.2 Clustering Jerárquico Aglomerativo (Agglomerative)}

\begin{itemize}
	\item \textbf{Mecanismo:} Construye una estructura en árbol (\emph{Dendrograma}) fusionando sucesivamente los dos clusters más cercanos desde un estado inicial donde cada punto es un cluster individual.
	
	\item \textbf{Criterios de Enlace (Linkage Criteria):}
	\begin{itemize}
		\item \textbf{Single Linkage:} Distancia mínima entre miembros de dos clusters ($d_{\min}(C_1, C_2)$).
		\item \textbf{Complete Linkage:} Distancia máxima entre miembros de dos clusters ($d_{\max}(C_1, C_2)$).
		\item \textbf{Average Linkage:} Distancia promedio entre todos los pares de elementos entre clusters.
	\end{itemize}
\end{itemize}

\subsection{6. Práctica y Aplicación en Código}

\begin{lstlisting}[language=Python, breaklines=true]
	import numpy as np
	import torch
	import torch.nn as nn
	from sklearn.neighbors import KNeighborsClassifier
	from sklearn.naive_bayes import GaussianNB
	from sklearn.tree import DecisionTreeClassifier
	from sklearn.cluster import KMeans
	from sklearn.preprocessing import StandardScaler
	
	# --------------------------------------------------------------------------
	# 1. Modelos Clasicos en Scikit-Learn (kNN, Naive Bayes, Decision Trees)
	# --------------------------------------------------------------------------
	def train_classical_models(X_train, y_train, X_test):
	# kNN: requiere estandarizacion previa de caracteristicas
	scaler = StandardScaler()
	X_train_scaled = scaler.fit_transform(X_train)
	X_test_scaled = scaler.transform(X_test)
	
	knn = KNeighborsClassifier(n_neighbors=5, metric='euclidean')
	knn.fit(X_train_scaled, y_train)
	
	# Gaussian Naive Bayes
	gnb = GaussianNB()
	gnb.fit(X_train, y_train)
	
	# Arbol de Decision CART (Criterio Gini)
	tree = DecisionTreeClassifier(criterion='gini', max_depth=5)
	tree.fit(X_train, y_train)
	
	return knn, gnb, tree
	
	# --------------------------------------------------------------------------
	# 2. Clustering (k-means)
	# --------------------------------------------------------------------------
	def run_kmeans_clustering(X_data, n_clusters=3):
	kmeans = KMeans(n_clusters=n_clusters, random_state=42, n_init=10)
	cluster_labels = kmeans.fit_predict(X_data)
	centroids = kmeans.cluster_centers_
	return cluster_labels, centroids
	
	# --------------------------------------------------------------------------
	# 3. MLP Baseline con Vectorizacion (Aplanado) en PyTorch
	# --------------------------------------------------------------------------
	class MLPClassifier(nn.Module):
	def __init__(self, input_shape=(3, 64, 64), num_classes=10, hidden_units=[512, 256], dropout_rate=0.3):
	super(MLPClassifier, self).__init__()
	self.input_dim = input_shape[0] * input_shape[1] * input_shape[2]
	
	layers = []
	in_dim = self.input_dim
	for h_dim in hidden_units:
	layers.append(nn.Linear(in_dim, h_dim))
	layers.append(nn.ReLU())
	layers.append(nn.Dropout(p=dropout_rate))
	in_dim = h_dim
	
	layers.append(nn.Linear(in_dim, num_classes))
	self.network = nn.Sequential(*layers)
	
	def forward(self, x):
	x = x.view(x.size(0), -1)
	return self.network(x)
\end{lstlisting}